\documentclass[a4paper]{article}

% Basic packages
% Español (México) con XeLaTeX: fontspec para fuentes Unicode, polyglossia
% para la división silábica y los nombres del documento en la variante mexicana.
\usepackage{fontspec}
\usepackage{polyglossia}
\setdefaultlanguage[variant=mexican]{spanish}
% Los nombres chinos van como «pinyin (original)»: xeCJK da a los caracteres
% Han su propia fuente; sin ella XeLaTeX los omite con un aviso.
% FandolSong no cubre algunos caracteres de nombres (U+73FA); WenQuanYi Zen Hei sí.
\usepackage[AutoFallBack=true]{xeCJK}
\setCJKmainfont{FandolSong-Regular.otf}
\setCJKfallbackfamilyfont{\CJKrmdefault}{WenQuanYi Zen Hei}
% polyglossia no traduce los nombres de listings.
\AtBeginDocument{\providecommand{\lstlistingname}{}\renewcommand{\lstlistingname}{Listado}%
  \providecommand{\lstlistlistingname}{}\renewcommand{\lstlistlistingname}{Índice de listados}}
\usepackage{amsmath, amssymb}
\usepackage{graphicx}
\usepackage[margin=2.5cm]{geometry}
\usepackage[most]{tcolorbox}
\usepackage{etoolbox}
\usepackage{listings}
\usepackage{hyperref}
\usepackage{booktabs}
\usepackage{subcaption}
\usepackage{float}
\usepackage{longtable}

% Highlight boxes
\newtcolorbox{knowledgebox}[1]{
    enhanced, colback=blue!5!white, colframe=blue!75!black, colbacktitle=blue!75!black,
    coltitle=white, fonttitle=\bfseries, title={#1}, attach boxed title to top left={yshift=-2mm, xshift=2mm},
    boxrule=1pt, sharp corners
}

\newtcolorbox{importantbox}[1]{
    enhanced, colback=yellow!10!white, colframe=yellow!80!black, colbacktitle=yellow!80!black,
    coltitle=black, fonttitle=\bfseries, title={#1}, sharp corners
}

\newtcolorbox{warningbox}[1]{
    enhanced, colback=red!5!white, colframe=red!75!black, colbacktitle=red!75!black,
    coltitle=white, fonttitle=\bfseries, title={#1}, sharp corners
}

% Listings
\lstset{
    language=Python,
    basicstyle=\ttfamily\small,
    keywordstyle=\color{blue},
    stringstyle=\color{red!60!black},
    commentstyle=\color{green!60!black},
    breaklines=true,
    frame=single,
    numbers=left,
    numberstyle=\tiny\color{gray},
    captionpos=b,
    extendedchars=true
}

% Front-page metadata
\newcommand{\notetitle}{{[}LLM Agents F24{]} Open-Source and Science in the Era of Foundation Models — Percy Liang}
\newcommand{\noteauthors}{Elaboradas a partir de materiales públicos del curso}
\newcommand{\notedate}{18 de noviembre de 2024}
\newcommand{\videochannel}{Berkeley RDI}
\newcommand{\videopublishdate}{18 de noviembre de 2024}
\newcommand{\videoduration}{aproximadamente 70 minutos}
\newcommand{\videourl}{https://www.youtube.com/live/f3KKx9LWntQ}
\newcommand{\videocoverpath}{cover.jpg}
\newcommand{\repourl}{https://github.com/hqhq1025/ai-course-notes}

% Footer with repo info
\usepackage{fancyhdr}
\pagestyle{fancy}
\fancyhf{}
\fancyfoot[C]{\thepage}
\fancyfoot[R]{\footnotesize\color{gray}\href{\repourl}{AI Course Notes}}
\renewcommand{\headrulewidth}{0pt}
\renewcommand{\footrulewidth}{0pt}

\begin{document}

\begin{titlepage}
\centering
{\Large Notas del curso\par}
\vspace{1.2cm}
{\huge\bfseries \notetitle\par}
\vspace{0.8cm}
{\large \noteauthors\par}
\vspace{0.3cm}
{\large \notedate\par}
\vspace{1.2cm}

\ifdefempty{\videocoverpath}{
{\small Al generar las notas, indique la ruta local de la portada del video.\par}
}{
\includegraphics[width=0.82\textwidth,height=0.45\textheight,keepaspectratio]{\videocoverpath}\par
}

\vfill
\begin{tcolorbox}[width=0.9\textwidth, colback=black!2!white, colframe=black!60, sharp corners]
\textbf{Autor o canal del video}: \videochannel\par
\textbf{Fecha de publicación}: \videopublishdate\par
\textbf{Duración del video}: \videoduration\par
\ifdefempty{\videourl}{}{
\textbf{Enlace del video}: \href{\videourl}{\nolinkurl{\videourl}}\par
}
\par
\textbf{Página del proyecto}: \href{\repourl}{\nolinkurl{\repourl}}\par
\end{tcolorbox}
\end{titlepage}

\tableofcontents
\newpage

%% --- Inicio del contenido --- %%

\section{Hilo conductor de la clase: la tijera entre el aumento de capacidades y la disminución de la apertura}

\subsection{Juicio inicial: Why Access Matters}
Percy Liang plantea, en los primeros minutos, la tesis central de toda la clase: \textbf{las capacidades de los Foundation Model aumentan rápidamente, pero el nivel de acceso disponible para los investigadores disminuye en paralelo}. Esta tijera de «curva de capacidad hacia arriba, curva de apertura hacia abajo» no es solo un cambio de estrategia comercial, sino que reconfigurará directamente la agenda de investigación de los próximos cinco años.

\begin{importantbox}{Access Shapes Research}
El expositor lo resume en una frase: \textit{``Access shapes research.''}
Cuando solo se tiene acceso por API, lo que más fácilmente se puede hacer es Prompt Engineering; cuando se cuenta con los pesos y el código de entrenamiento, se abre la posibilidad de hacer interpretabilidad mecanicista, modificación del entrenamiento, auditoría de seguridad e innovación arquitectónica.
\end{importantbox}

\begin{figure}[H]
\centering
\includegraphics[width=0.93\textwidth]{frame-01-title.jpg}
\caption{Página de apertura de la clase: del tema de los Agent hacia una pregunta abierta de ciencia}
\end{figure}
\noindent\footnotesize Evidencia visual: los subtítulos se ubican aproximadamente entre 00:00:09--00:00:51, donde el expositor plantea explícitamente la ``importance of open source'' y ``access shapes research''.\normalsize

\subsection{Perspectiva histórica: los tres “dividendos de acceso” del NLP}
El expositor repasa varios saltos característicos en la historia del NLP y el aprendizaje automático:
\begin{enumerate}
    \item de los años noventa a los dosmiles: la mayor accesibilidad de texto en internet impulsó el NLP estadístico;
    \item la década de 2010: la maduración de las plataformas de etiquetado de datos impulsó el aprendizaje supervisado y la evaluación mediante benchmarks;
    \item la era del aprendizaje profundo: la mayor disponibilidad de GPU y clústeres de cómputo hizo posible escalar el transformer.
\end{enumerate}

\begin{knowledgebox}{Metodología implícita a nivel de curso}
Este repaso histórico no es un «paréntesis de divulgación», sino que prepara una metodología para lo que sigue:
\textbf{al discutir las capacidades de un modelo, es necesario discutir a la vez el acceso a los datos, el acceso al cómputo y el acceso al sistema.}
\end{knowledgebox}

\begin{warningbox}{Error común: reducir la pregunta de apertura a «si los pesos son de código abierto»}
La pregunta de apertura no es un interruptor binario. Que los pesos puedan descargarse es solo una dimensión. Si los datos de entrenamiento, el proceso de entrenamiento, los scripts de evaluación, los registros del sistema y la estrategia de despliegue no son visibles, la reproducibilidad de la investigación sigue siendo insuficiente y las conclusiones científicas siguen siendo frágiles.
\end{warningbox}

\subsection{Resumen de la sección}
Esta sección plantea la pregunta central de toda la clase: no se trata de debatir si «el código abierto es bueno o no», sino de discutir qué investigación permite y qué investigación bloquea cada nivel de acceso, y cómo esta diferencia estructural afecta el progreso a largo plazo de los Agent y los Foundation Model.

\section{Tres niveles de acceso: API, Open Weights, Open Source}

\subsection{El modelo de tres niveles de acceso}
El expositor divide el acceso en tres niveles, cada uno correspondiente a una posición de investigación distinta:
\begin{itemize}
    \item API access: como un «científico conductual» que observa las entradas y salidas de una caja negra;
    \item Open weights: como un «neurocientífico» que puede examinar las representaciones internas;
    \item Open source: como un «ingeniero de sistemas» que puede reescribir todo el pipeline.
\end{itemize}

\begin{table}[H]
\centering
\begin{tabular}{p{2.4cm}p{4.2cm}p{4.8cm}}
\toprule
\textbf{Nivel de acceso} & \textbf{Capacidades típicas} & \textbf{Limitaciones principales} \\
\midrule
API access & Diseño de prompts, orquestación de llamadas a herramientas, evaluación de caja negra & Difícil verificar las causas mecánicas, difícil hacer modificaciones a nivel estructural \\
Open weights & Análisis de representaciones, destilación/fine-tuning, experimentos de intervención & Aun así, sigue restringido por la receta de entrenamiento y la arquitectura existentes \\
Open source & Reconstrucción de toda la cadena de datos, modelo, entrenamiento y evaluación & Exige mucho de la ingeniería y del cómputo, el costo de reproducción es alto \\
\bottomrule
\end{tabular}
\caption{El marco de «niveles de acceso» propuesto por el expositor y sus fronteras de investigación}
\end{table}

\begin{importantbox}{Un mismo modelo, distintos niveles de acceso llevan a distintos tipos de artículo}
El recordatorio más valioso del curso es este: la pregunta de investigación no está determinada únicamente por el interés, sino también por «qué experimento es posible hacer».
En la era de la API son más frecuentes los artículos de estrategia o de nivel de prompt; en la era de open source son más probables los artículos de arquitectura, objetivo de entrenamiento, construcción de datos y codiseño de sistemas.
\end{importantbox}

\subsection{Por qué hay que distinguir Open Weights de Open Source}
\begin{knowledgebox}{Open Weights no equivale a Open Science}
Tener los pesos permite hacer muchas cosas, pero aun así no se sabe:
\begin{itemize}
    \item cómo influye la proporción de mezcla de datos en las fronteras de las capacidades;
    \item cómo cambia el comportamiento el currículo de entrenamiento (curriculum) y las estrategias de alineación;
    \item si la mejora de desempeño proviene de la arquitectura o de la ingeniería de datos.
\end{itemize}
Por eso, Open Weights es un estado intermedio importante, pero todavía no llega al estado científico de «poder verificarse por completo».
\end{knowledgebox}

\begin{warningbox}{Riesgo de confusión conceptual en las discusiones de política}
En las discusiones de la industria suele mezclarse «poder invocar la API», «poder descargar los pesos» y «poder reproducir el entrenamiento» en una sola palabra: «abierto». Esto hace que el diseño de la regulación y de los mecanismos de colaboración se aleje de las restricciones técnicas reales y genere incentivos equivocados.
\end{warningbox}

\subsection{Resumen de la sección}
El marco de tres niveles que propone el expositor permite convertir el «debate sobre la apertura» de un eslogan en un problema de ingeniería: distintos niveles corresponden a distintas capacidades experimentales y distintas exposiciones al riesgo, y no deben confundirse entre sí.

\section{Agentes en la era de las API: de función universal a sistema ejecutable}

\subsection{La perspectiva de la función universal}
En condiciones de solo-API, el LLM puede considerarse como una «función universal condicionada por lenguaje natural». Tiene una gran generalidad a nivel textual, por lo que el peso de la ingeniería se traslada a «cómo incrustar esta función en un flujo ejecutable».

\begin{importantbox}{La API no es un nodo hoja, sino un controlador}
El docente enfatiza: en los sistemas de agentes modernos, la API del LLM no es solo una subllamada dentro de un programa fijo; con frecuencia determina directamente el flujo de ejecución, la selección de herramientas y la actualización de estados intermedios, desempeñando en realidad el papel de controlador.
\end{importantbox}

\begin{figure}[H]
\centering
\includegraphics[width=0.93\textwidth]{frame-03-abstraction.jpg}
\caption{Agentic Abstraction: abstracción unificada de modelo, herramientas y humano}
\end{figure}
\noindent\footnotesize Evidencia visual: subtítulos aproximadamente en 00:04:18--00:05:30, el docente comenta «API controls execution flow» y «memory stream / retrieval / action».\normalsize

\subsection{Flujo de memoria, recuperación y ciclo de reflexión}
La arquitectura básica de agente que presenta el docente puede descomponerse en cuatro componentes de un ciclo cerrado: observación (observation), flujo de memoria (memory stream), recuperación (retrieval) y ejecución (action).
Además, el sistema incorpora un mecanismo de «reflexión» (reflection) que hace que el modelo produzca resúmenes de alto nivel de las trayectorias pasadas, y luego reincorpora esos resúmenes a las decisiones posteriores.

\begin{knowledgebox}{Por qué la reflexión es tan importante en los agentes}
El fallo de un agente basado en llamadas a herramientas rara vez proviene de un único error de inferencia, sino de un sesgo acumulado a lo largo de una trayectoria de varios pasos.
El valor de la reflexión está en comprimir el «ruido de la trayectoria» en una señal de estrategia reutilizable, de modo que los pasos posteriores dejen de repetirse a ciegas.
\end{knowledgebox}

\begin{warningbox}{Las dos trampas más frecuentes de los sistemas de memoria}
\begin{itemize}
    \item Acumulación sin límite: el contexto se expande de manera continua y degrada la calidad de la recuperación;
    \item Fijación de errores: los resúmenes erróneos se citan de manera repetida y forman una «alucinación autorreforzada».
\end{itemize}
Por eso, la memoria no consiste en «almacenar cuanto más mejor», sino en que importa más la calidad de la compresión y la capacidad de verificación.
\end{warningbox}

\subsection{Resumen de la sección}
Lo esencial del agente en la era de las API no es «llamar a unas cuantas herramientas más», sino construir un ciclo de ejecución controlable: hacer del LLM un controlador y, al mismo tiempo, usar mecanismos de memoria y reflexión para reducir el error acumulado en cadenas largas.

\section{Orquestación multi-Agent: de la libertad estructural a la disciplina de ingeniería}

\subsection{Dimensiones de la orquestación}
El curso ofrece dimensiones de diseño claras para la orquestación multi-Agent: estático/dinámico, centralizado/descentralizado, contexto compartido/aislado, automatizado/con intervención humana. Estas dimensiones determinan el límite superior del desempeño del sistema y el límite inferior de su estabilidad.

\begin{figure}[H]
\centering
\includegraphics[width=0.93\textwidth]{frame-04-orchestration.jpg}
\caption{Dimensiones de diseño de Multi-Agent Orchestration}
\end{figure}
\noindent\footnotesize Evidencia visual: los subtítulos se ubican aproximadamente en 00:06:02--00:07:40, cuando el orador aborda la resolución de tareas complejas y la organización de trayectorias.\normalsize

\begin{importantbox}{Multi-Agent no es «múltiples ventanas en paralelo», sino «diseño de interfaces de tarea»}
Si no hay límites de rol, límites de contexto y límites de aceptación, aumentar el número de Agents por lo regular solo amplifica los conflictos y el trabajo duplicado.  
El paralelismo verdaderamente eficaz proviene de contratos claros de entrada y salida, así como de formatos claros para los productos intermedios.
\end{importantbox}

\subsection{Patrones de diseño: Conversation / Tool Use / Planning / Memory}
\begin{figure}[H]
\centering
\includegraphics[width=0.93\textwidth]{frame-05-patterns.jpg}
\caption{Agentic Design Patterns: colaboración entre conversación, herramientas, planificación y memoria}
\end{figure}
\noindent\footnotesize Evidencia visual: los subtítulos se ubican aproximadamente en 00:07:50--00:09:10, cuando el orador analiza sucesivamente patrones como reflection, plan y fact-check.\normalsize

\begin{knowledgebox}{Cómo los cuatro tipos de patrones forman un sistema mínimo viable}
\begin{itemize}
    \item Conversation: sostiene la colaboración entre roles y la transmisión de contexto;
    \item Tool Use: traduce las decisiones lingüísticas en acciones ejecutables;
    \item Planning: descompone tareas largas en subobjetivos verificables;
    \item Memory: conserva y comprime el estado clave a través de los turnos.
\end{itemize}
\end{knowledgebox}

\begin{warningbox}{Limitarse a «apilar patrones» hace que la complejidad del sistema se salga de control}
Los patrones no son mejores cuantos más se usen. Cada mecanismo que se agrega debe responder dos preguntas:  
¿qué tipo de tasa de fallos reduce? ¿qué tipo de nuevo modo de falla introduce?
\end{warningbox}

\subsection{Resumen de la sección}
El núcleo de la ingeniería multi-Agent no es la «orquestación creativa», sino el «diseño de restricciones». Solo cuando las interfaces son claras, el estado es verificable y la reversión es ejecutable, la orquestación compleja puede implementarse de manera estable.

\section{Problem-Solving Agents: el caso de MLAgentBench}

\subsection{Convertir «saber conversar» en «saber hacer»}
El expositor presentó una dirección representativa de los problem-solving agents: completar un ciclo de extremo a extremo en tareas reales de aprendizaje automático, que incluye leer la tarea, modificar el código, ejecutar experimentos, analizar registros e iterar la estrategia.

\begin{figure}[H]
\centering
\includegraphics[width=0.93\textwidth]{frame-06-groupchat.jpg}
\caption{Planeación y resolución de tareas complejas: de la conversación a un pipeline ejecutable}
\end{figure}
\noindent\footnotesize Evidencia en pantalla: los subtítulos se ubican aproximadamente en 00:10:02--00:13:10, donde el expositor discute la resolución de tareas del benchmark, la ejecución en bash y la evaluación automática.\normalsize

\begin{importantbox}{Valor del benchmark: convertir la «discusión sobre capacidades» en un «experimento reproducible»}
Sin un benchmark, la mejora del agente suele quedarse en la narrativa de una demostración.  
Con un benchmark, es posible comparar la ganancia real entre distintos modelos, distintas estrategias y distintas cadenas de herramientas, y descubrir la distribución de los fallos.
\end{importantbox}

\subsection{Objetivos de evaluación y descomposición de fallos}
\begin{table}[H]
\centering
\begin{tabular}{p{2.8cm}p{4.0cm}p{4.8cm}}
\toprule
\textbf{Dimensión de evaluación} & \textbf{Señal observable} & \textbf{Falla típica} \\
\midrule
Grado de cumplimiento de la tarea & Si la métrica alcanza el objetivo, si el envío es válido & Desviación del objetivo, malentendido de la tarea \\
Confiabilidad de la ejecución & Tasa de éxito de los comandos, número de reintentos, tasa de tiempos de espera agotados & Errores de dependencia del entorno, fallas en la llamada a herramientas \\
Calidad de la estrategia & Eficiencia de la iteración, cobertura de la búsqueda, eficacia de la reflexión & Prueba y error a ciegas, óptimo local \\
Seguridad y restricciones & Comportamiento fuera de los permisos otorgados, efectos secundarios anómalos, capacidad de auditoría & Operación no autorizada, imposibilidad de revertir \\
\bottomrule
\end{tabular}
\caption{La evaluación de un problem-solving agent no solo mira la «puntuación final», sino también la calidad de la trayectoria de ejecución}
\end{table}

\begin{knowledgebox}{Por qué las cyber tasks se convierten en un escenario de prueba clave}
Las tareas de ciberseguridad tienen de manera natural una alta densidad de retroalimentación (éxito o fracaso claramente definidos) y un alto riesgo operativo (el costo de una operación errónea es alto).  
Este tipo de tareas permite examinar al mismo tiempo la capacidad de ejecución del agente, su capacidad de seguir restricciones y su capacidad de recuperación.
\end{knowledgebox}

\begin{warningbox}{«Inflar la puntuación» no equivale a una mejora real de la capacidad}
Si la distribución del benchmark es demasiado estrecha, el agente puede aprender rutinas específicas en lugar de una estrategia general.  
El expositor enfatiza que la distribución de tareas y las dimensiones de evaluación deben ampliarse de manera continua, para evitar un progreso aparente de «puntuación alta pero no transferible».
\end{warningbox}

\subsection{Resumen de la sección}
El avance clave de los problem-solving agent proviene de ser «ejecutable + evaluable + revisable». Lo que trabajos del tipo MLAgentBench impulsan realmente es un paradigma de investigación, y no solo el resultado de una sola tabla de clasificación.

\section{Simulation Agents: del razonamiento sobre tareas al modelado de procesos sociales}

\subsection{Por qué hacer simulation}
El ponente pasa después a los simulation agents: no solo hacer que el agente resuelva tareas, sino que también interactúe de manera continua en un contexto social, para estudiar el comportamiento colectivo, la difusión de información y los mecanismos de colaboración.

\begin{importantbox}{El valor de simulation no está en «parecerse a los humanos», sino en «poder formar hipótesis verificables»}
Cuando el agente participa como individuo programable en un sistema multiagente, podemos cambiar de manera sistemática la estrategia de memoria, las reglas de comunicación y la función de incentivos, y observar cómo cambia el comportamiento macroscópico.
\end{importantbox}

\subsection{Los tres niveles de mecanismo de la sociedad generativa}
\begin{knowledgebox}{El puente de lo microscópico a lo macroscópico}
\begin{enumerate}
    \item Nivel de memoria: cómo el individuo conserva y recupera sus experiencias;
    \item Nivel de decisión: cómo el individuo actúa dentro de un contexto;
    \item Nivel de interacción: cómo los individuos se influyen entre sí y forman una dinámica colectiva.
\end{enumerate}
Solo cuando estos tres niveles operan de manera conjunta se constituye un «sistema de simulación social susceptible de investigarse».
\end{knowledgebox}

\begin{warningbox}{Riesgo de validez externa de los resultados de la simulación}
Las reglas en un sistema de simulación suelen ser fijadas por quien investiga, de manera que los resultados pueden reflejar más la «elección de reglas» que un mecanismo social real.
Sin una validación suficiente frente a la realidad, las conclusiones de simulation deben considerarse generación de hipótesis, y no demostración de hechos.
\end{warningbox}

\subsection{Resumen de la sección}
Los simulation agents extienden la investigación de Agent al «nivel del comportamiento del sistema». Su producto más valioso no son las historias, sino hipótesis de mecanismo que pueden someterse a prueba experimental y ser refutadas.

\section{Ventana de investigación de Open-Weight: interpretación, ataque y defensa}

\subsection{Oportunidades de investigación que trae la visibilidad interna}
Al entrar en la capa open-weight, las cosas que los investigadores pueden hacer aumentan de forma significativa: la interpretabilidad de representaciones, el análisis de activaciones, la edición de características, la destilación y los experimentos de alineación se vuelven posibles.

\begin{importantbox}{Lo que dan los Open Weights es «derecho de investigación mecanicista», no «derecho de modificación universal»}
Puedes profundizar en el interior del modelo, pero sigues estando limitado por el historial de entrenamiento, la forma de la arquitectura y el origen de los datos.  
Por lo tanto, open-weight es una capa de gran capacidad, pero no es la capa final.
\end{importantbox}

\subsection{Lado de la seguridad: la mejora de capacidad y la exposición al riesgo ocurren al mismo tiempo}
En la parte media de la clase se mencionan varias veces experimentos relacionados con la seguridad: los pesos abiertos permiten que la investigación de defensa sea más profunda, y también hacen que las rutas de ataque sean más viables.  
Esto no es una pregunta de opción única de «abrir o no abrir», sino una pregunta sistémica de «cómo acompañarlo con mecanismos de gobernanza y auditoría».

\begin{table}[H]
\centering
\begin{tabular}{p{2.8cm}p{4.0cm}p{4.8cm}}
\toprule
\textbf{Dirección de investigación} & \textbf{Beneficio de los pesos abiertos} & \textbf{Riesgo correspondiente} \\
\midrule
Interpretación mecanicista & Permite localizar patrones de falla y características vulnerables & Los atacantes también pueden localizar los puntos débiles \\
Experimentos de alineación & Permite reproducir el entrenamiento y las rutas de intervención de seguridad & Replicar variantes dañinas a bajo costo \\
Evaluación de defensa & Permite realizar pruebas de equipo rojo y pruebas de estrés de alta intensidad & Los detalles de la defensa son aprovechados por el aprendizaje adversarial \\
\bottomrule
\end{tabular}
\caption{Estructura de «aumento simultáneo de capacidad y riesgo» en la investigación de open-weight}
\end{table}

\begin{warningbox}{Un error de juicio común en la comunidad técnica: externalizar la gobernanza de riesgos al «proveedor del modelo»}
Una vez que el ecosistema entra en una forma de código abierto o de pesos abiertos, la gobernanza de riesgos no puede depender de una sola organización.  
Se necesita que los puntos de referencia, las herramientas de auditoría, las normas de despliegue y los mecanismos de responsabilidad evolucionen en conjunto.
\end{warningbox}

\subsection{Resumen de la sección}
Open-weight abrió una ventana importante para la investigación de mecanismos y la investigación de seguridad, pero al mismo tiempo amplificó la dificultad de la gobernanza del ecosistema. El avance técnico debe expandirse en conjunto con la capacidad de gobernanza.


\section{El significado estricto de Open Source: de «utilizable» a «reproducible»}

\subsection{La naturaleza técnica de la disputa sobre la definición}
En la segunda mitad, el ponente discutió el problema de los límites de la definición de apertura: si solo se cuenta con los pesos y faltan los datos de entrenamiento y el proceso completo, ¿basta eso para llamarlo open source? No se trata de una disputa semántica, sino de algo que afecta la reproducibilidad de la investigación científica y la delimitación de responsabilidades.

\begin{importantbox}{Open Source for Models necesita un nuevo marco de definición}
El open source tradicional del software gira principalmente en torno a que «el código fuente sea visible y modificable».  
En los sistemas de modelos, los datos, el proceso de entrenamiento, el protocolo de evaluación y la estrategia de publicación son igualmente componentes de la «fuente».
\end{importantbox}

\subsection{La escalera de la reproducibilidad}
\begin{table}[H]
\centering
\begin{tabular}{p{3.0cm}p{3.8cm}p{4.8cm}}
\toprule
\textbf{Nivel de apertura} & \textbf{Reproducción posible} & \textbf{Lo que aún falta} \\
\midrule
API disponible & Reproducción del comportamiento, evaluación a nivel de interfaz & Atribución del mecanismo, rastreo del entrenamiento \\
Weights disponibles & Reproducción de la inferencia, reproducción del fine-tuning, experimentos parciales de interpretación & La cadena causal de datos y entrenamiento \\
Full source & Reproducción de extremo a extremo del entrenamiento y la evaluación, modificación a nivel de sistema & El principal cuello de botella pasa a ser el cómputo y los recursos de ingeniería \\
\bottomrule
\end{tabular}
\caption{La escalera de capacidades de «poder invocarse» a «poder reproducirse»}
\end{table}

\begin{knowledgebox}{Por qué «reproducible» es más importante que «descargable»}
La comunidad científica necesita una acumulación de conocimiento verificable. Solo cuando las condiciones experimentales clave son reproducibles, comparables y refutables, el campo puede formar un consenso sólido, en lugar de depender de las afirmaciones no verificables de organizaciones individuales.
\end{knowledgebox}

\begin{warningbox}{Tratar «open source» como una etiqueta de mercadotecnia termina por dañar la confianza de la comunidad}
Si el término se usa mal de manera sostenida, los investigadores cometerán errores sistemáticos en sus expectativas de colaboración:  
creerán que algo es reproducible cuando en realidad no lo es; creerán que algo admite auditoría cuando en realidad no la admite.
\end{warningbox}

\subsection{Resumen de la sección}
El núcleo de esta sección es que «la definición es infraestructura». Para los sistemas de modelos, la definición de apertura determina directamente la reproducibilidad de la investigación, la eficiencia de la colaboración y la operatividad de la gobernanza.


\section{Cómputo y restricciones organizacionales: el cuello de botella real de la ciencia abierta}

\subsection{El cómputo sigue siendo la restricción de primer orden}
El ponente señala que, incluso con una definición clara de apertura y una comunidad con voluntad firme, la ciencia abierta seguirá estando limitada por el cómputo y por la ingeniería de sistemas.
Esto significa que la competencia futura no solo estará en la capacidad de los modelos, sino también en quién pueda construir infraestructura de entrenamiento y evaluación reutilizable.

\begin{figure}[H]
\centering
\includegraphics[width=0.93\textwidth]{frame-07-adoption.jpg}
\caption{Adopción del ecosistema y contribución de la comunidad: los efectos de red de los marcos abiertos}
\end{figure}
\noindent\footnotesize Evidencia en pantalla: subtítulos aproximadamente en 00:26:01 y 00:31:01, donde el ponente aborda la escala del ecosistema, las prácticas de los modelos abiertos y el valor de la migración de la investigación.\normalsize

\begin{knowledgebox}{Tres tipos de recursos escasos desde la perspectiva de la infraestructura}
\begin{itemize}
    \item Cómputo de entrenamiento: determina el espacio de modelos y datos que se puede explorar;
    \item Talento de ingeniería: determina la mantenibilidad del sistema y el rendimiento de los experimentos;
    \item Activos de evaluación: determinan si el progreso se puede verificar y comparar.
\end{itemize}
\end{knowledgebox}

\subsection{Una vía viable para la colaboración descentralizada}
\begin{importantbox}{La verdadera palanca del ecosistema abierto: estándares y cadenas de herramientas compartidos}
El cómputo de un solo punto difícilmente puede replicar el entrenamiento frontier, pero la comunidad puede formar un «multiplicador de colaboración» mediante evaluaciones unificadas, scripts de entrenamiento portables y procesos reutilizables de gobernanza de datos.
\end{importantbox}

\begin{warningbox}{Sin estandarización, la colaboración degenera en trabajo repetido y fragmentado}
Si cada equipo usa formatos de registro, protocolos de evaluación y definiciones de tareas incompatibles entre sí, el resultado es que «todos están haciendo experimentos, pero nadie puede compararlos de manera efectiva».
\end{warningbox}

\subsection{Resumen de la sección}
La ciencia abierta no equivale a una barrera de entrada baja. Requiere nueva infraestructura común para reducir la fricción de la colaboración; de lo contrario, la «apertura» se quedará solo en el nivel del lema.
\section{Lista de acciones para equipos de investigación}

\subsection{Del curso a la ingeniería: cómo empezar a construir un flujo de desarrollo de Agentes reproducible}
\begin{figure}[H]
\centering
\includegraphics[width=0.93\textwidth]{frame-08-research-questions.jpg}
\caption{Cierre de las preguntas de investigación: el triángulo de capacidad, seguridad y participación humana}
\end{figure}
\noindent\footnotesize Evidencia visual: subtítulos aproximadamente en 00:51:01--00:52:00, el docente enfatiza los temas sistémicos y las líneas de investigación posteriores.\normalsize

\begin{importantbox}{Ruta mínima recomendada para la implementación}
\begin{enumerate}
    \item Definir primero tareas verificables, y discutir después la arquitectura del Agente;
    \item Establecer primero registros y reproducción, y optimizar después la política;
    \item Hacer primero la descomposición de fallos, y buscar después el puntaje en las tablas de clasificación;
    \item Unificar primero el protocolo de evaluación, y ampliar después el trabajo paralelo del equipo.
\end{enumerate}
\end{importantbox}

\subsection{Panel recomendado de evaluación y gobernanza}
\begin{table}[H]
\centering
\begin{tabular}{p{2.8cm}p{4.1cm}p{4.6cm}}
\toprule
\textbf{Dimensión del panel} & \textbf{Métricas clave} & \textbf{Recomendación de implementación} \\
\midrule
Desempeño en la tarea & pass rate, success@k, latency & Fijar la versión de la tarea y conservar trayectorias históricas reproducibles \\
Calidad de ejecución & tasa de fallos de herramientas, tasa de reintentos, tasa de reversiones & Introducir registros de auditoría obligatorios y etiquetas de clasificación de fallos \\
Restricciones de seguridad & eventos de exceso de privilegios, alertas de operaciones sensibles & Delimitar con claridad los límites de permisos y los puntos de confirmación humana \\
Reproducibilidad científica & varianza de los resultados, consistencia del entorno, número de repeticiones del experimento & Publicar instantáneas de configuración y versiones de los scripts de evaluación \\
\bottomrule
\end{tabular}
\caption{Panel de monitoreo para que la investigación de Agentes pase de «funcionar» a «acumularse»}
\end{table}

\begin{knowledgebox}{Un diseño organizacional clave en la colaboración de equipos}
Quien dirige la investigación debe establecer la reproducibilidad como objetivo de primer orden:
los commits de código, las versiones de datos, las versiones de modelo y las versiones de evaluación deben poder formar una misma cadena de trazabilidad, para evitar «un resultado correcto pero un proceso imposible de rastrear».
\end{knowledgebox}

\begin{warningbox}{No busques la «automatización total» de manera prematura}
En tareas de alto riesgo y alta incertidumbre, una intervención humana razonable no es un retroceso, sino un mecanismo necesario para aumentar la robustez del sistema.
El verdadero objetivo es la «automatización controlada», no la «fantasía de operar sin supervisión».
\end{warningbox}

\subsection{Resumen de la sección}
En el plano de la implementación, lo más importante no es el nombre de algún marco de trabajo, sino si estas cuatro cosas pueden sostenerse a la vez: la definición de la tarea, el protocolo de evaluación, la gobernanza de registros y el control de permisos.


\section{Hoja de ruta con casos: sprint de 12 semanas de investigación abierta en Agent}

\subsection{De la filosofía del curso al plan de ejecución}
Para convertir la filosofía de esta clase en un proyecto ejecutable, se puede adoptar una estructura de «sprint de 12 semanas»: las primeras 4 semanas se dedican a construir la infraestructura, las 4 semanas intermedias a la iteración del sistema, y las últimas 4 semanas a la robustez y al empaquetado para la reproducción. Este ritmo corresponde al espíritu central del expositor: \textbf{primero convertir el sistema de investigación en una máquina verificable, y solo después buscar picos de desempeño locales.}

\begin{figure}[H]
\centering
\includegraphics[width=0.93\textwidth]{frame-02-benefits.jpg}
\caption{Objetivo de beneficios del sistema Agent: la trinidad de interfaz, capacidad y arquitectura}
\end{figure}
\noindent\footnotesize Evidencia visual: los subtítulos se ubican aproximadamente entre 00:00:43 y 00:01:39; el expositor extiende «access shapes research» hacia los límites de capacidad del sistema.\normalsize

\begin{importantbox}{El objetivo sugerido del sprint no es «lanzar una demo llamativa»}
Un objetivo de mayor valor es:
\begin{itemize}
    \item que cualquier conclusión experimental pueda ser reproducida por un compañero en 24 horas;
    \item que cualquier trayectoria de fallo pueda ubicarse rápidamente en una etapa concreta;
    \item que cualquier mejora pueda explicarse como proveniente de los datos, la estrategia o la cadena de herramientas.
\end{itemize}
Hacer esto parece lento en el corto plazo, pero a largo plazo reduce de manera significativa la proporción de «progreso falso».
\end{importantbox}

\subsection{Hitos por etapa y entregables}
\begin{table}[H]
\centering
\begin{tabular}{p{2.2cm}p{3.5cm}p{3.6cm}p{3.2cm}}
\toprule
\textbf{Etapa} & \textbf{Tarea central} & \textbf{Entregable obligatorio} & \textbf{Criterio de aceptación} \\
\midrule
Semanas 1--2 & Unificar el protocolo de tareas y el formato de registro & task schema, run schema, replay script & Cualquier tarea puede reproducirse \\
Semanas 3--4 & Hacer funcionar el agent de línea base y la evaluación & baseline result, failure taxonomy & Existe una varianza de línea base estable \\
Semanas 5--6 & Incorporar los módulos de memoria, reflexión y planificación & ablation report, trace dashboard & Se puede explicar el origen de la ganancia \\
Semanas 7--8 & Orquestación de múltiples agent y gobernanza de permisos & orchestration spec, safety gate & Las conductas que exceden los permisos pueden bloquearse y auditarse \\
Semanas 9--10 & Robustez y pruebas de estrés fuera de distribución & stress suite, adversarial suite & Los fallos bajo distintas distribuciones pueden descomponerse \\
Semanas 11--12 & Empaquetado para la reproducción y publicación documentada & reproducibility package, model card & Terceros pueden reproducirlo de manera independiente \\
\bottomrule
\end{tabular}
\caption{Ruta mínima de ingeniería de investigación del sprint de 12 semanas}
\end{table}

\begin{knowledgebox}{Por qué hacer primero el replay y después la optimización}
Muchos equipos dedican el 80\% del tiempo a «pensar nuevas estrategias» sin contar con trayectorias reproducibles.  
Sin replay, no se puede determinar si una mejora se debe a que la estrategia funciona o a una fluctuación aleatoria; sin un replay estable, ninguna optimización posterior puede acumularse como conocimiento confiable.
\end{knowledgebox}

\subsection{Ciclo principal experimental sugerido (pseudocódigo)}
\begin{lstlisting}[caption={Research Loop: ciclo principal de iteración reproducible de Agent},label={lst:research-loop}]
for task_batch in benchmark:
    run_id = launch_agent(task_batch, config, tools, safety_policy)
    traces = collect_traces(run_id)
    metrics = evaluate(traces, eval_protocol)

    if metrics.regression_detected:
        rollback(config)
        tag_failure(traces, taxonomy)
        continue

    insights = reflection_and_ablation(traces, metrics)
    config = update_config(config, insights)
    snapshot(run_id, config, metrics, traces)
\end{lstlisting}

\begin{warningbox}{La deuda técnica más fácil de pasar por alto: experimentos no comparables}
Si en cada iteración se cambian el conjunto de tareas, la versión de las herramientas o el script de evaluación, surge la ilusión de que «siempre se está progresando, pero no hay dos resultados que puedan compararse».  
Durante la etapa del sprint es necesario congelar los protocolos clave y ajustar las variables solo dentro de una ventana controlada.
\end{warningbox}

\subsection{Resumen de la sección}
El objetivo de esta hoja de ruta de 12 semanas es convertir la filosofía de «ciencia abierta + ingeniería de Agent» en un sistema ejecutable: reproducible, auditable e interpretable. Si falta cualquiera de los tres, el equipo tendrá que rehacer el trabajo repetidamente en la etapa de escalamiento.


\section{Referencia rápida de términos y métricas}

\subsection{Tabla de correspondencia de términos clave}
\begin{longtable}{p{2.7cm}p{3.2cm}p{5.6cm}}
\toprule
\textbf{Término} & \textbf{Comprensión en español} & \textbf{Función en esta clase} \\
\midrule
\endfirsthead
\toprule
\textbf{Término} & \textbf{Comprensión en español} & \textbf{Función en esta clase} \\
\midrule
\endhead
API Access & Acceso de caja negra a la interfaz & Permite experimentación a nivel de comportamiento, pero dificulta la atribución a nivel de mecanismo. \\
Open Weights & Pesos abiertos & Permite el análisis interno, el fine-tuning y la destilación, pero sigue restringido por la receta de entrenamiento existente. \\
Open Source & Apertura de toda la cadena & Abarca datos, modelo, entrenamiento y evaluación, y permite la reproducibilidad a nivel de sistema. \\
Access Shapes Research & El acceso determina la forma de la investigación & Hilo conductor del curso: los recursos accesibles determinan qué tipo de preguntas se pueden plantear. \\
Agent Controller & Controlador de ejecución & El LLM no es solo un módulo al que se invoca, sino el núcleo de control del flujo. \\
Memory Stream & Flujo de memoria & Registra observaciones y acciones históricas, y da contexto para la recuperación y la reflexión posteriores. \\
Retrieval & Recuperación & Extrae del historial de memoria el estado más relevante para la tarea actual. \\
Reflection & Reflexión & Comprime la trayectoria histórica en una estrategia de alto nivel, lo que reduce la repetición de errores. \\
Planning & Planificación & Descompone una tarea larga en subobjetivos verificables. \\
Tool Use & Uso de herramientas & Convierte las decisiones del lenguaje en acciones de ejecución reales. \\
Orchestration & Orquestación & Organiza los límites de roles, los límites de contexto y el orden de ejecución de múltiples agentes. \\
Benchmark & Evaluación de referencia & Convierte la impresión de «parece más fuerte» en un resultado comparable y verificable. \\
Failure Taxonomy & Taxonomía de fallas & Distingue tipos de falla como malentendido del objetivo, falla de herramienta o degradación de la estrategia. \\
Replayability & Capacidad de repetición & Permite que cada falla se pueda repetir y diagnosticar; es la base de la reproducibilidad. \\
Audit Log & Registro de auditoría & Aporta la cadena de evidencia para la seguridad y la rendición de cuentas. \\
Distribution Shift & Cambio de distribución & Evalúa la estabilidad y la capacidad de generalización del agente en escenarios fuera de distribución. \\
Safety Gate & Compuerta de seguridad & Establece confirmación de permisos y mecanismos de reversión para acciones de alto riesgo. \\
Reproducibility Package & Paquete de reproducibilidad & Reúne configuración, scripts, entorno y resultados para que un tercero los verifique. \\
Open Definition & Definición de apertura & Precisa los límites de «apertura» para evitar que el uso confuso del término distorsione la colaboración. \\
Research Throughput & Rendimiento de la investigación & Lo determinan conjuntamente el cómputo, la cadena de herramientas y los procesos organizacionales. \\
\bottomrule
\end{longtable}

\begin{knowledgebox}{Por qué hacer una referencia rápida de términos}
La investigación de agentes abarca modelos, sistemas, evaluación y gobernanza, y el vocabulario que usan distintos equipos no siempre coincide.
Estandarizar los términos puede reducir de manera significativa los malentendidos en la colaboración, sobre todo cuando se reproducen experimentos entre instituciones.
\end{knowledgebox}

\subsection{Resumen de la sección}
Unificar la terminología es una acción de ingeniería de bajo costo y alto rendimiento. No mejora directamente el puntaje en las tablas de clasificación, pero sí mejora de manera significativa la comparabilidad de los experimentos del equipo y la velocidad de acumulación de conocimiento.

\section{Síntesis y ampliación}

\subsection{Tabla de síntesis de toda la clase}
\begin{table}[H]
\centering
\begin{tabular}{p{2.8cm}p{4.1cm}p{4.6cm}}
\toprule
\textbf{Tema} & \textbf{Idea central del conferencista} & \textbf{Implicación para la investigación en Agentes} \\
\midrule
Tendencia de apertura & El aumento de capacidad y la disminución del acceso ocurren simultáneamente & Los límites de la investigación quedan moldeados directamente por los permisos de acceso \\
Marco de tres niveles de acceso & La estratificación en API / Open Weights / Open Source es clara & Cada nivel corresponde a una capacidad distinta de experimentación verificable \\
Ingeniería de Agentes basados en API & El LLM funge como controlador de ejecución & El énfasis se desplaza hacia la memoria, la reflexión y la gestión de trayectorias \\
Benchmark y evidencia empírica & Contar con un sistema evaluable es condición previa para el progreso & Se pasa de una narrativa basada en demos a comparaciones reproducibles \\
Definición de apertura y gobernanza & Los límites terminológicos influyen en las políticas y la colaboración & «Descargable» no equivale a «reproducible» \\
Ruta a largo plazo & La capacidad técnica y la infraestructura coevolucionan & La ciencia abierta requiere la coordinación de estándares, cadenas de herramientas y cómputo \\
\bottomrule
\end{tabular}
\caption{Síntesis de los argumentos centrales de Lecture 03}
\end{table}

\subsection{Lecturas adicionales}
\begin{itemize}
    \item Bommasani et al., \textit{On the Opportunities and Risks of Foundation Models}, 2021.
    \item Liang et al., \textit{Holistic Evaluation of Language Models (HELM)}, 2022.
    \item Artículos de referencia relacionados con Agent evaluation and reproducibility (incluida la serie de trabajos de MLAgentBench).
    \item Discusiones en curso y actualizaciones de borradores de la OSI sobre la definición de apertura de los modelos de IA.
    \item Prácticas de ingeniería en marcos de trabajo con múltiples Agentes: comparación de diseño entre AutoGen, LangGraph y CrewAI.
\end{itemize}

\subsection{Idea de cierre}
\textit{``Access shapes research.''} Esta frase, aparentemente simple, en realidad define los límites reales de los investigadores en la era de los Agentes.  
Si se desea lograr un progreso científico sostenible en las direcciones de Foundation Model y Agent, la comunidad debe avanzar simultáneamente en la capacidad de los modelos, el acceso abierto y la infraestructura reproducible, en lugar de perseguir únicamente un SOTA puntual.

%% --- Fin del contenido --- %%

\end{document}
